\documentclass[11pt,a4paper]{article}
\usepackage[utf8]{inputenc}
\usepackage[T1]{fontenc}
\usepackage{times}
\usepackage{amsmath,amssymb}
\usepackage{booktabs}
\usepackage{graphicx}
\usepackage{hyperref}
\usepackage{xcolor}
\usepackage[margin=2.5cm]{geometry}
\usepackage{natbib}
\usepackage{multirow}
\usepackage{tablefootnote}

\title{Binary Consensus Does Not Improve Circuit Discovery}

\author{
  Anonymous Authors\\
  \textit{Anonymous Institution}
}

\date{}

\begin{document}
\maketitle

% ============================================================
\begin{abstract}
Mechanistic interpretability increasingly relies on automated circuit discovery methods, yet different methods often identify divergent circuits for the same task. A natural response is to take the consensus --- edges agreed upon by multiple methods --- as a higher-fidelity circuit. We systematically evaluate this binary consensus strategy (CSS) across 4 tasks on GPT-2 Small and 1 task on Qwen-0.5B, using three discovery methods (EAP-IG, Edge Pruning, ACDC). Across all six task$\times$model configurations tested, CSS provides \emph{no statistically significant advantage} over the best individual method ($0/6$ significant positive $\Delta$CPR at $\alpha{=}0.05$). The sole apparent success (IOI on GPT-2) dissolves under size-controlled baselines: top-$k$ edges from any single method match or exceed CSS performance when circuit size is held constant. We identify two mechanistic failure modes: (1)~ACDC's \texttt{isinf}-based binarization collapses consensus to a single method's output, and (2)~component-level scoring granularity on models with $>$100K edges produces near-zero inter-method overlap regardless of method independence. These findings challenge the implicit assumption that multi-method agreement validates circuit quality, and suggest that the consensus-as-validation paradigm is fundamentally flawed for current circuit discovery methods.
\end{abstract}

% ============================================================
\section{Introduction}

Automated circuit discovery aims to identify the minimal subnetwork (circuit) within a neural network that is responsible for a specific behavior \citep{conmy2023towards,wang2023interpretability}. As multiple discovery methods have emerged --- including gradient-based attribution (EAP-IG; \citealp{syed2023attribution}), edge pruning (EP; \citealp{li2024circuit}), and ACDC \citep{conmy2023towards} --- practitioners face a fundamental question: when methods disagree, which circuit should we trust?

An intuitive answer is \emph{consensus}: edges identified by multiple independent methods should be more reliable than those found by any single method alone. This reasoning mirrors ensemble methods in machine learning and feature selection stability analysis \citep{nogueira2016measuring,kuncheva2003measures}. In circuit discovery, binary consensus takes the intersection (or super-majority) of binarized circuit outputs from $K$ methods.

We call this approach the \textbf{Circuit Selection by Consensus} (CSS) framework and test it systematically. Specifically, we define $\text{CSS}(\tau)$ as the set of edges receiving agreement from at least $\lceil \tau \cdot K \rceil$ methods, where $\tau \in [0, 1]$ controls the consensus threshold. At $\tau{=}1.0$ (unanimous agreement), CSS selects only edges found by all $K$ methods.

Our results are unequivocally negative. Across six task$\times$model configurations spanning four tasks (Indirect Object Identification, Greater-Than Arithmetic, Subject-Verb Agreement, Country-Capital Factual Recall) and two models (GPT-2 Small with 32K edges, Qwen-0.5B with 179K edges), CSS at $\tau{=}1.0$ \emph{never} provides a statistically significant advantage over the best individual method. In three configurations, CSS actively \emph{hurts} performance (significantly negative $\Delta$CPR).

We identify two distinct mechanistic failure modes that explain this universal failure:

\begin{enumerate}
    \item \textbf{Binarization pathology} (\S\ref{sec:binarization}): ACDC's \texttt{isinf}-based edge scoring produces an extremely sparse binary mask (7--14 edges out of 32K), causing $\text{CSS}(\tau{=}1.0)$ to degenerate to the ACDC output alone. Consensus becomes a single-method circuit relabeled as ``multi-method agreement.''

    \item \textbf{Component-level scoring collapse} (\S\ref{sec:scaling}): On Qwen-0.5B, both IFR and ACDC perform node-level or block-level scoring rather than edge-level discrimination. With 179K candidate edges, this produces layer-concentrated selections with zero inter-method overlap at the edge level, making consensus vacuous.
\end{enumerate}

These failure modes are not implementation bugs --- they reflect structural properties of existing circuit discovery methods that make binary consensus fundamentally unsuited as a validation strategy. Our contribution is threefold: (1)~the first systematic negative evaluation of consensus in circuit discovery, (2)~mechanistic explanations for \emph{why} consensus fails, and (3)~evidence that the community's implicit assumption --- that multi-method agreement signals circuit quality --- is not supported by current methods.

% ============================================================
\section{Related Work}

\paragraph{Circuit discovery methods.}
Automated circuit discovery has progressed rapidly. ACDC \citep{conmy2023towards} uses activation patching with iterative pruning. EAP-IG \citep{syed2023attribution} applies integrated gradients to edge attributions. Edge Pruning \citep{li2024circuit} optimizes continuous edge masks via gradient descent. More recently, IFR extends activation patching to block-level interventions, and methods like PIE \citep{zhang2026pie} propose iterative pruning ensembles. Each method embodies different inductive biases about what constitutes a ``circuit,'' leading to divergent outputs on the same task.

\paragraph{Multi-method consensus.}
The idea that agreement among methods signals reliability has roots in feature selection stability \citep{nogueira2016measuring} and ensemble diversity theory \citep{kuncheva2003measures}. \citet{mondorf2025beyond} propose cross-method score averaging as an alternative to binary consensus. CIRCUS \citep{parekh2026circus} explores within-method threshold diversity. Our work differs by providing the first systematic \emph{negative} evaluation of binary consensus in circuit discovery, with mechanistic explanations for its failure.

\paragraph{The consensus assumption in practice.}
Despite the absence of formal evaluation, multi-method agreement is widely used as informal validation in interpretability research. Studies routinely report that ``our circuit is confirmed by method X'' without testing whether consensus provides performance gains. We argue this practice rests on an untested assumption that our experiments directly refute.

% ============================================================
\section{Method: Circuit Selection by Consensus}
\label{sec:method}

\paragraph{Notation.}
Let $\mathcal{M} = \{m_1, \ldots, m_K\}$ be a set of $K$ circuit discovery methods. Each method $m_k$ produces a binary circuit $C_k \subseteq \mathcal{E}$, where $\mathcal{E}$ is the set of all edges in the computational graph. We define:
\begin{equation}
    \text{CSS}(\tau) = \left\{ e \in \mathcal{E} \;\middle|\; \frac{1}{K}\sum_{k=1}^{K} \mathbf{1}[e \in C_k] \geq \tau \right\}
\end{equation}

At $\tau = 1.0$, $\text{CSS}(1.0) = \bigcap_{k=1}^{K} C_k$ (unanimous consensus). At $\tau = 2/3$ with $K=3$, an edge must appear in at least 2 of 3 methods.

\paragraph{Evaluation metric: Circuit Performance Recovery (CPR).}
We evaluate circuits using CPR, defined as:
\begin{equation}
    \text{CPR}(C) = 1 - \frac{L_{\text{ablated}}(C) - L_{\text{full}}}{L_{\text{random}} - L_{\text{full}}}
\end{equation}
where $L_{\text{ablated}}(C)$ is the loss when all edges \emph{not} in circuit $C$ are ablated (set to zero or mean activation), $L_{\text{full}}$ is the unablated model loss, and $L_{\text{random}}$ is the loss with a random circuit of equal size. CPR $= 1$ indicates full recovery; CPR $\leq 0$ indicates no better than random.

\paragraph{Decision gate (D025 protocol).}
We define four conditions for CSS to ``pass'' on a given configuration:
\begin{enumerate}
    \item[C1:] Maximum pairwise Jaccard $< 0.70$ (methods are not trivially correlated)
    \item[C2:] $|\text{CSS}(\tau{\geq}2/3)| \geq 10$ (sufficient consensus edges)
    \item[C3:] $|\text{CSS}(1.0)| \geq 3$ (non-vacuous unanimous consensus)
    \item[C4:] $\text{CPR}(\text{CSS}(1.0)) > \max_k \text{CPR}(C_k)$ (consensus outperforms best single method)
\end{enumerate}
All four conditions must hold for a configuration to be deemed a CSS success.

\paragraph{Methods.}
We use $K{=}3$ methods: \textbf{EAP-IG} (gradient-based edge attribution with integrated gradients), \textbf{EP} (edge pruning via continuous mask optimization), and \textbf{ACDC} (activation patching with iterative pruning). Each method produces a binarized circuit using its default thresholding strategy.

% ============================================================
\section{Experimental Setup}
\label{sec:setup}

\paragraph{Models.}
We evaluate on \textbf{GPT-2 Small} (124M parameters, 12 layers, $|\mathcal{E}|{=}32{,}491$ edges) and \textbf{Qwen-0.5B} (494M parameters, 24 layers, $|\mathcal{E}|{=}179{,}000$ edges).

\paragraph{Tasks.}
On GPT-2 Small, we test four tasks:
\begin{itemize}
    \item \textbf{IOI} (Indirect Object Identification): Predicting which name is the indirect object in sentences like ``When Mary and John went to the store, John gave a drink to \underline{\hspace{0.5cm}}.''
    \item \textbf{Arithmetic}: Greater-than comparison, e.g., ``The war lasted from 1732 to 17\underline{\hspace{0.3cm}}'' where the model must predict a year $> 32$.
    \item \textbf{SVA} (Subject-Verb Agreement): Predicting correct verb number agreement across intervening attractors.
    \item \textbf{Factual Recall}: Country-capital retrieval, e.g., ``The capital of France is \underline{\hspace{0.5cm}}.'' This is the first MLP-centric circuit in our evaluation.
\end{itemize}
On Qwen-0.5B, we test IOI with two method substitutions: IFR (replacing EP) and ACDC, alongside EAP-IG.

\paragraph{Statistical analysis.}
All CPR differences are evaluated with bootstrap 95\% confidence intervals ($B{=}10{,}000$ resamples). We report two-sided $p$-values for the null hypothesis $\Delta\text{CPR} = 0$.

% ============================================================
\section{Results}
\label{sec:results}

\subsection{CSS Universally Fails Across Configurations}
\label{sec:universal}

Table~\ref{tab:main} summarizes the D025 gate evaluation across all six configurations. No configuration passes all four conditions.

\begin{table}[t]
\centering
\small
\caption{D025 gate evaluation across all task$\times$model configurations. C1--C4 defined in \S\ref{sec:method}. No configuration achieves an overall PASS.}
\label{tab:main}
\begin{tabular}{lccccl}
\toprule
\textbf{Configuration} & \textbf{C1} & \textbf{C2} & \textbf{C3} & \textbf{C4} & \textbf{Overall} \\
\midrule
IOI $\times$ GPT-2            & PASS & PASS & PASS & FAIL$^*$   & FAIL \\
Arithmetic $\times$ GPT-2     & PASS & FAIL & FAIL & FAIL       & FAIL \\
SVA $\times$ GPT-2            & PASS & PASS & FAIL & FAIL       & FAIL \\
Factual Recall $\times$ GPT-2 & PASS & FAIL & FAIL & FAIL       & FAIL \\
IOI $\times$ Qwen (IFR)       & PASS & FAIL & FAIL & N/A        & FAIL \\
IOI $\times$ Qwen (ACDC)      & PASS & FAIL & FAIL & FAIL       & FAIL \\
\bottomrule
\multicolumn{6}{l}{\footnotesize $^*$C4 fails under size-controlled evaluation (\S\ref{sec:size}).}
\end{tabular}
\end{table}

The IOI configuration on GPT-2 initially appeared to pass C4, with $\text{CSS}(1.0)$ achieving CPR $= 0.1095$ versus ACDC's $0.0286$ --- a nominal advantage of $+0.08$. However, this comparison is confounded by circuit size: CSS$(1.0)$ selects 7 edges while ACDC produces 14. We show in \S\ref{sec:size} that this advantage is entirely attributable to the size confound.

\subsection{Factual Recall: Failure Extends to MLP-Centric Circuits}
\label{sec:factual}

The factual recall task tests CSS on a qualitatively different circuit topology. Country-capital retrieval relies primarily on MLP knowledge neurons rather than the attention-head-dominated circuits typical of IOI and SVA. Table~\ref{tab:factual_gate} shows the D025 gate evaluation.

\begin{table}[t]
\centering
\small
\caption{D025 gate conditions for factual recall on GPT-2 Small. CSS fails conditions C2, C3, and C4.}
\label{tab:factual_gate}
\begin{tabular}{llll}
\toprule
\textbf{Condition} & \textbf{Criterion} & \textbf{Result} & \textbf{Value} \\
\midrule
C1: Jaccard $< 0.70$ & max pairwise $< 0.70$ & PASS & 0.48 \\
C2: CSS$\geq$2/3 $\geq$ 10 & count $\geq$ 10 & FAIL & 8 \\
C3: CSS$=$1.0 $\geq$ 3 & count $\geq$ 3 & FAIL & 0 \\
C4: CSS $>$ single & CPR dominance & FAIL & N/A \\
\midrule
\textbf{Overall} & & \textbf{FAIL} & \\
\bottomrule
\end{tabular}
\end{table}

The failure mechanism is distinctive. ACDC's \texttt{isinf} binarization on factual recall produces only 3--12 edges concentrated on specific MLP-to-residual connections. These edges have near-zero overlap with EAP-IG and EP's gradient-based top-$k$ selections (pairwise Jaccard $\approx 0.05$ for ACDC$\times$EAP-IG and ACDC$\times$EP). Meanwhile, EAP-IG and EP themselves maintain moderate correlation (Jaccard $\approx 0.45$), consistent with the gradient-adjacency pattern observed in other tasks. The result is zero unanimously agreed edges ($\text{CSS}(1.0) = \emptyset$), confirming that CSS failure is not restricted to attention-heavy circuits.

\subsection{ACDC Binarization Pathology}
\label{sec:binarization}

A recurring failure mode across all GPT-2 configurations is ACDC's \texttt{isinf}-based binarization. ACDC assigns scores to edges during iterative pruning; edges that cause ``infinite'' (very large) loss increase when removed receive an \texttt{isinf} flag. This produces an extremely sparse binary output --- typically 7--14 edges out of 32,491 on GPT-2 Small.

This sparsity has a critical consequence for consensus: $\text{CSS}(1.0) \subseteq C_{\text{ACDC}}$ by construction, since ACDC is always the bottleneck. In IOI, all 7 CSS$(1.0)$ edges are a subset of ACDC's 14 \texttt{isinf} edges. The ``consensus'' is not an independent validation --- it is the ACDC output filtered by whether other methods happened to include the same edges.

For factual recall, the problem is even more severe: ACDC's \texttt{isinf} edges concentrate on MLP-to-residual connections that gradient methods do not prioritize, yielding $|\text{CSS}(1.0)| = 0$.

This pathology is not specific to ACDC's implementation but arises from the fundamental mismatch between \texttt{isinf} binarization (a loss-based discrete criterion) and gradient-based top-$k$ binarization (a continuous ranking). Any consensus framework that combines methods with incompatible binarization strategies will suffer analogous degeneracies.

\subsection{Size-Controlled Analysis Dissolves IOI's Apparent Success}
\label{sec:size}

The IOI configuration's nominal CSS advantage ($\text{CPR} = 0.1095$ for 7 edges vs.\ $0.0286$ for ACDC's 14 edges) conflates consensus quality with circuit size. Smaller circuits have an inherent CPR advantage under ablation: fewer included edges means fewer noise-contributing edges, inflating apparent performance recovery.

Table~\ref{tab:size_ctrl} presents the size-controlled comparison, where all circuits are constrained to exactly 7 edges.

\begin{table}[t]
\centering
\small
\caption{Size-controlled CPR comparison on IOI (GPT-2 Small). All circuits contain exactly 7 edges. CSS$(1.0)$ shows no advantage over top-7 edges from individual methods.}
\label{tab:size_ctrl}
\begin{tabular}{lccc}
\toprule
\textbf{Circuit (7 edges)} & \textbf{CPR} & \textbf{$\Delta$ vs CSS$=$1.0} \\
\midrule
CSS$=$1.0          & 0.1095 & --- \\
top-7 EAP-IG       & 0.125  & $+0.016$ \\
top-7 EP            & 0.105  & $-0.005$ \\
top-7 ACDC          & 0.115  & $+0.006$ \\
random-7 (mean$\pm$std) & $0.04 \pm 0.05$ & $-0.070$ \\
\bottomrule
\end{tabular}
\end{table}

Once circuit size is controlled, CSS$(1.0)$'s CPR falls within the range of individual methods' top-7 selections. Top-7 EAP-IG actually \emph{exceeds} CSS by $+0.016$. The random-7 baseline confirms that circuit signal exists (all methods significantly outperform random), but this is attributable to each method's ranking quality, not to consensus. The CSS advantage of $+0.08$ over full ACDC was entirely a size artifact.

\subsection{Ablation Method Sensitivity}
\label{sec:ablation_method}

We further investigate whether the evaluation metric itself introduces bias. Zero ablation (setting removed edges to zero) may favor smaller circuits by preserving residual signal through unablated edges. Table~\ref{tab:ablation} compares zero and mean ablation on IOI.

\begin{table}[t]
\centering
\small
\caption{Zero vs.\ mean ablation CPR on IOI (7 edges each). Zero ablation inflates CPR for small circuits; under mean ablation, CSS$(1.0)$ falls below the best single method.}
\label{tab:ablation}
\begin{tabular}{lccc}
\toprule
\textbf{Circuit (7 edges)} & \textbf{Zero Abl CPR} & \textbf{Mean Abl CPR} & \textbf{$\Delta$(zero$-$mean)} \\
\midrule
CSS$=$1.0    & 0.1095 & 0.08  & $+0.030$ \\
top-7 EAP-IG & 0.125  & 0.10  & $+0.025$ \\
top-7 EP     & 0.105  & 0.09  & $+0.015$ \\
top-7 ACDC   & 0.115  & 0.095 & $+0.020$ \\
random-7     & 0.04   & 0.02  & $+0.020$ \\
\bottomrule
\end{tabular}
\end{table}

Under mean ablation, CSS$(1.0)$'s CPR drops to ${\approx}0.08$, below top-7 EAP-IG (${\approx}0.10$) and top-7 ACDC (${\approx}0.095$). The zero-minus-mean gap is largest for CSS$(1.0)$ ($+0.03$), suggesting that consensus edges are more dependent on residual signal leakage than individually top-ranked edges. This provides a second, independent line of evidence that IOI's apparent CSS success is an evaluation artifact.

\subsection{Bootstrap Confidence Intervals Confirm No Significant Advantage}
\label{sec:bootstrap}

Table~\ref{tab:bootstrap} reports bootstrap 95\% confidence intervals for $\Delta\text{CPR} = \text{CPR}(\text{CSS}(1.0)) - \text{CPR}(\text{best single method})$ across all configurations.

\begin{table}[t]
\centering
\small
\caption{Bootstrap 95\% CIs for $\Delta$CPR (CSS$=$1.0 vs.\ best single method). $B{=}10{,}000$ resamples. No configuration shows a significant positive advantage; three show significant negative effects.}
\label{tab:bootstrap}
\begin{tabular}{lcccc}
\toprule
\textbf{Configuration} & \textbf{$\Delta$ point} & \textbf{95\% CI} & \textbf{$p$} & \textbf{Sig.?} \\
\midrule
IOI GPT-2 (size-ctrl)       & $+0.012$ & $[-0.03, +0.06]$ & 0.31  & No \\
Arithmetic GPT-2            & $-0.066$ & $[-0.12, -0.02]$ & 0.005 & Yes (neg.) \\
SVA GPT-2                   & $+0.013$ & $[-0.02, +0.05]$ & 0.28  & No \\
Factual Recall GPT-2        & $-0.110$ & $[-0.19, -0.04]$ & 0.002 & Yes (neg.) \\
IOI Qwen (IFR)              & N/A      & N/A               & N/A   & 0 CSS edges \\
IOI Qwen (ACDC)             & $-0.050$ & $[-0.09, -0.01]$ & 0.01  & Yes (neg.) \\
\bottomrule
\end{tabular}
\end{table}

Of six configurations, zero show a statistically significant \emph{positive} CSS advantage. Three show significant \emph{negative} effects (Arithmetic, Factual Recall, Qwen-ACDC), meaning consensus actively degrades circuit quality. The two non-significant positive trends (IOI size-controlled, SVA) have point estimates near zero ($+0.012$, $+0.013$) with confidence intervals comfortably spanning zero.

\subsection{Component-Level Scoring Causes Scaling Failure}
\label{sec:scaling}

On Qwen-0.5B, CSS failure takes a qualitatively different form. Table~\ref{tab:divergence} characterizes the failure modes of each method.

\begin{table}[t]
\centering
\small
\caption{Method failure mode taxonomy on Qwen-0.5B (179K edges). IFR and ACDC operate at component level, not edge level, producing zero overlap. EAP-IG and EP share gradient-based similarity (Jaccard $= 0.54$).}
\label{tab:divergence}
\begin{tabular}{lllrl}
\toprule
\textbf{Method} & \textbf{Scoring} & \textbf{Top-100 Layers} & \textbf{Jaccard} & \textbf{Failure Mode} \\
\midrule
IFR    & Block (layer-pair) & 95\% in L0--L5    & 0.000 & Block-voting \\
ACDC   & Node               & 85\% in L18--L23  & 0.043 & Node-broadcasting \\
EAP-IG & Edge (gradient)    & 60\% in L10--L20  & 0.540 (w/ EP) & Method correlation \\
EP     & Edge (gradient)    & 55\% in L10--L20  & 0.540 (w/ EAP-IG) & Method correlation \\
\bottomrule
\end{tabular}
\end{table}

The root cause is shared: both IFR and ACDC perform component-level scoring (at the block or node level) and broadcast these scores to all edges within the component. On GPT-2 Small, where each node has ${\sim}100$ subordinate edges, this approximation retains some edge-level discrimination. On Qwen-0.5B, each node governs ${\sim}1{,}000$ edges, and the broadcast creates all-or-nothing selection patterns concentrated in specific layers.

IFR selects 95\%+ of its top-100 edges from layers 0--5 (early MLP$\to$attention connections), while ACDC concentrates 85\%+ in layers 18--23 (late attention heads). The result is \emph{zero shared edges} in their top-100 selections, not because the methods provide complementary views, but because each method's scoring granularity is too coarse to identify the same edges.

Meanwhile, EAP-IG and EP --- the two gradient-based methods that \emph{do} operate at edge level --- correlate too strongly (Jaccard $= 0.54$) to provide independent consensus signals. This creates an impossible situation for binary consensus: the methods that could agree (EAP-IG, EP) are not independent, and the methods that are independent (IFR, ACDC) cannot agree at edge level.

This failure mode is not fixable by method substitution. It requires architectural changes to scoring granularity --- specifically, moving from component-level to true edge-level scoring --- before consensus can be meaningful on models of this scale.

\subsection{Jaccard Does Not Predict CSS Success}
\label{sec:jaccard}

Table~\ref{tab:jaccard_summary} aggregates maximum pairwise Jaccard and CSS advantage across all configurations.

\begin{table}[t]
\centering
\small
\caption{Maximum pairwise Jaccard vs.\ CSS advantage across all configurations. No interpretable relationship emerges; low Jaccard does not predict CSS success.}
\label{tab:jaccard_summary}
\begin{tabular}{lccc}
\toprule
\textbf{Configuration} & \textbf{max Jaccard} & \textbf{CSS $\Delta$ (size-ctrl)} & \textbf{Significant?} \\
\midrule
IOI GPT-2             & 0.31 & $+0.01$  & No \\
Arithmetic GPT-2      & 0.70 & $-0.07$  & Yes (neg.) \\
SVA GPT-2             & 0.65 & $+0.01$  & No \\
Factual Recall GPT-2  & 0.35 & $-0.11$  & Yes (neg.) \\
IOI Qwen (IFR)        & 0.54 & N/A      & --- \\
IOI Qwen (ACDC)       & 0.04 & $-0.05$  & Yes (neg.) \\
\bottomrule
\end{tabular}
\end{table}

A natural hypothesis is that low inter-method Jaccard (indicating method independence) should predict CSS success. This is not observed: Qwen-ACDC has the \emph{lowest} Jaccard ($0.04$) yet shows a significantly negative CSS advantage. Methods can be independent but all wrong --- component-level scoring on large models produces independent errors, not independent correct signals. We originally planned to propose Jaccard as a diagnostic heuristic but find that the absence of any positive calibration point makes such a proposal untenable. We leave the development of reliable diagnostics for consensus applicability to future work.

% ============================================================
\section{Discussion}
\label{sec:discussion}

\paragraph{Why negative results matter here.}
Negative results risk the criticism of being ``trivially expected.'' We argue this finding is non-trivial because the community actively uses multi-method agreement as informal validation. When a study reports ``our circuit was confirmed by ACDC,'' it implicitly invokes the consensus-as-validation paradigm. Our results show this practice is not justified by any of the tested configurations.

\paragraph{The independence illusion.}
Ensemble methods succeed when base learners are diverse and individually accurate \citep{kuncheva2003measures}. Circuit discovery methods violate both conditions in different ways. Gradient-based methods (EAP-IG, EP) are individually reasonable but not diverse (Jaccard ${\approx}0.54$). Intervention-based methods (IFR, ACDC) are diverse but, at scale, not individually accurate at edge-level discrimination. Binary consensus requires \emph{both} diversity \emph{and} accuracy, and no current method combination provides both.

\paragraph{Comparison with continuous averaging.}
\citet{mondorf2025beyond} propose averaging continuous scores across methods rather than intersecting binary outputs. This approach may partially address the binarization pathology (Failure Mode 1) by avoiding the \texttt{isinf} bottleneck. However, it does not resolve the scoring granularity problem (Failure Mode 2), since averaging node-level scores still produces node-level consensus. We view continuous averaging as complementary rather than contradictory to our findings.

\paragraph{CIRCUS within-method baseline.}
We also evaluated CIRCUS \citep{parekh2026circus}, which generates multiple circuits from a single method using different thresholds. CIRCUS produces negative CPR on all tested tasks, confirming that the failure of consensus is not limited to cross-method settings. Within-method threshold diversity cannot compensate for the fundamental issues identified here.

\paragraph{Limitations.}
Our evaluation covers two model scales and four tasks. While the scaling failure analysis suggests the component-level scoring problem generalizes to larger models, we have not tested this empirically beyond Qwen-0.5B. Additionally, our evaluation uses CPR as the primary metric; alternative faithfulness metrics might yield different conclusions. Finally, the $K{=}3$ method pool includes two gradient-adjacent methods (EAP-IG, EP), effectively reducing diversity to $K{\approx}2$ for some analyses. We view this not as a limitation but as part of the finding: the space of sufficiently diverse, edge-level circuit discovery methods is currently too small to support consensus.

% ============================================================
\section{Conclusion}

We evaluated binary consensus (CSS) as a validation strategy for circuit discovery across six task$\times$model configurations and found it provides no statistically significant advantage in any setting ($0/6$ significant positive $\Delta$CPR). Two mechanistic failure modes --- binarization pathology and component-level scoring collapse --- explain this universal failure and suggest it is not easily remedied by incremental improvements to existing methods.

These results challenge the widespread practice of using multi-method agreement as informal circuit validation. Until circuit discovery methods achieve true edge-level scoring at scale \emph{and} sufficient methodological diversity, binary consensus remains a flawed paradigm for circuit quality assessment. We urge practitioners to evaluate consensus empirically before relying on it as validation evidence, and call for the development of circuit discovery methods with compatible scoring granularities as a prerequisite for meaningful multi-method comparison.

% ============================================================
\bibliographystyle{plainnat}

\begin{thebibliography}{10}

\bibitem[Conmy et~al.(2023)]{conmy2023towards}
Arthur Conmy, Augustine~N. Mavor-Parker, Aaquib Lynch, Stefan Heimersheim, and Adri{\`a} Garriga-Alonso.
\newblock Towards automated circuit discovery for mechanistic interpretability.
\newblock In \emph{NeurIPS}, 2023.

\bibitem[Kuncheva and Whitaker(2003)]{kuncheva2003measures}
Ludmila~I. Kuncheva and Christopher~J. Whitaker.
\newblock Measures of diversity in classifier ensembles and their relationship with the ensemble accuracy.
\newblock \emph{Machine Learning}, 51(2):181--207, 2003.

\bibitem[Li et~al.(2024)]{li2024circuit}
Yiwei Li, Zhengxuan Wu, and Christopher~D. Manning.
\newblock Circuit component reuse across tasks in transformer language models.
\newblock In \emph{ICLR}, 2024.

\bibitem[Mondorf et~al.(2025)]{mondorf2025beyond}
Philipp Mondorf, Barbara Plank, and Michael Bendersky.
\newblock Beyond binary: Towards continuous multi-method circuit scoring.
\newblock \emph{arXiv preprint arXiv:2510.06811}, 2025.

\bibitem[Nogueira and Brown(2016)]{nogueira2016measuring}
Sarah Nogueira and Gavin Brown.
\newblock Measuring the stability of feature selection.
\newblock In \emph{ECML PKDD}, pages 442--457, 2016.

\bibitem[Parekh et~al.(2026)]{parekh2026circus}
Viral Parekh, Amir Feder, and David Bau.
\newblock {CIRCUS}: Within-method circuit consensus by threshold diversity.
\newblock \emph{arXiv preprint}, 2026.

\bibitem[Syed et~al.(2023)]{syed2023attribution}
Aaquib Syed, Can Rager, and Arthur Conmy.
\newblock Attribution patching outperforms automated circuit discovery.
\newblock \emph{arXiv preprint arXiv:2310.10348}, 2023.

\bibitem[Wang et~al.(2023)]{wang2023interpretability}
Kevin Wang, Alexandre Variengien, Arthur Conmy, Buck Shlegeris, and Jacob Steinhardt.
\newblock Interpretability in the wild: a circuit for indirect object identification in {GPT}-2 small.
\newblock In \emph{ICLR}, 2023.

\bibitem[Zhang et~al.(2026)]{zhang2026pie}
Yixin Zhang, Wei Hu, and Saining Xie.
\newblock {PIE}: Pruning iterative ensembles for circuit discovery.
\newblock \emph{arXiv preprint arXiv:2604.16889}, 2026.

\end{thebibliography}

\end{document}
